\documentclass[11pt]{article}
\usepackage[margin=1in]{geometry}
\usepackage{amsmath,amssymb}
\usepackage{microtype}
\usepackage[hidelinks]{hyperref}

\title{A smooth-symbol partial answer to the\\
spectral open-disk question in arXiv:2005.09557}
\author{Literature-implied identification note}
\date{17 August 2026}

\begin{document}
\maketitle

\begin{center}
\fbox{\parbox{0.9\textwidth}{\textbf{Status:}
literature-implied answer (partial smooth-symbol subcase).  No full answer is
claimed for an arbitrary $H^\infty$ analytic part.}}
\end{center}

\section*{Original question}

Abakumov, Baranov, Charpentier, and Lishanskii consider on $H^2$ the Toeplitz
operator with meromorphic symbol
\[
 \Phi(z)=R(1/z)+\varphi(z),\qquad \varphi\in H^\infty,
\]
where $R$ has degree $N$ and has no pole in the closed disk.  They prove that
hypercyclicity forces $N$-valence and record the spectrum formula
\[
 \sigma(T_\Phi)=\mathbb C\setminus\Phi(\mathbb D,N).
\]
They then ask whether
\[
 T_\Phi\ \text{hypercyclic}
 \quad\Longrightarrow\quad
 \sigma(T_\Phi)\cap\mathbb D\ne\varnothing.
\]
The question occurs on PDF page 3 of arXiv:2005.09557, immediately after the
spectrum formula \cite{ABCL}.

\section*{Supporting theorem and implication}

On PDF page 14, Theorem 3.2 of Fricain, Grivaux, and Ostermann \cite{FGO}
says that if a Toeplitz symbol $F$ satisfies

\begin{itemize}
\item[(H1)] $F\in C^{1+\varepsilon}(\mathbb T)$ for a suitable
$\varepsilon>0$, and $F'$ never vanishes;
\item[(H2)] $F(\mathbb T)$ has only finitely many simple self-intersections;
\item[(H3)] $\operatorname{wind}_F(\lambda)\le0$ off $F(\mathbb T)$;
\end{itemize}

then hypercyclicity of $T_F$ forces every connected component of
$\operatorname{int}\sigma(T_F)$ to meet the unit circle.

This gives a positive answer to the source question under (H1)--(H2).
Indeed, the source's Proposition 1.1 says that hypercyclicity makes $\Phi$
$N$-valent.  For $\lambda\notin\Phi(\mathbb T)$, the argument principle gives
\[
 \operatorname{wind}_\Phi(\lambda)=n_\Phi(\lambda)-N\le0,
\]
so (H3) is automatic.  Moreover, the regular immersed curve
$\Phi(\mathbb T)$ has a bounded complementary face adjacent along a regular
arc to the unbounded face.  The winding number is zero on the unbounded face
and changes by one across that arc; nonpositivity forces the bounded face to
have winding $-1$.  It is therefore a nonempty component of
$\operatorname{int}\sigma(T_\Phi)$.

By Theorem 3.2 this open component contains a point of $\mathbb T$.  Every
neighborhood of a unit-circle point contains points of $\mathbb D$, hence
\[
 \sigma(T_\Phi)\cap\mathbb D\ne\varnothing.
\]

\section*{Provenance and limitations}

The supporting authors cite the original paper and devote a comparison
section to its meromorphic-symbol setting, but the implication above is an
agent-identified combination rather than a named answer in their paper.
The unrestricted $\varphi\in H^\infty$ question is not resolved: the proof of
the supporting theorem uses a bounded $H^\infty$ functional calculus coming
from smooth Toeplitz model theory.  Radial smoothing is insufficient because
hypercyclicity is not stable under strong or norm perturbation.

Bounded searches used the exact source sentence, both arXiv identifiers,
the paper titles, and variants of ``hypercyclic Toeplitz spectrum open unit
disk.''  No full later answer was found.

\begin{thebibliography}{9}
\bibitem{ABCL}
E. Abakumov, A. Baranov, S. Charpentier, and A. Lishanskii,
\emph{New classes of hypercyclic Toeplitz operators},
arXiv:2005.09557 (2020; revised 2021).

\bibitem{FGO}
E. Fricain, S. Grivaux, and M. Ostermann,
\emph{Hypercyclicity of Toeplitz operators with smooth symbols},
arXiv:2502.03303, Theorem 3.2.
\end{thebibliography}

\end{document}
